\section{Variantes de Normalización y Control de Logits de Atención}
\label{sec:annex-normalization}

Este anexo documenta las variantes de normalización implementadas por el conjunto de herramientas y los mecanismos de normalización query--key que controlan el crecimiento de los logits de atención. Las secciones principales remiten a este anexo para las formulaciones y el pseudocódigo; las claves bibliográficas se resuelven contra \texttt{docs/bibliography/}.

\subsection{Variantes de Normalización: LayerNorm, RMSNorm, pRMSNorm, Tanh Dinámico, Erf Dinámico y FlashNorm}

La normalización determina cómo se controla la escala de activación a través de la profundidad. En este repositorio, la normalización no es solo una elección de modelado sino también una cuestión de compatibilidad con el esquema, porque solo ciertos valores son actualmente aceptados por \texttt{norm\_type}. El contrato del esquema es:
\[
\texttt{norm\_type} \in \{\texttt{layer\_norm}, \texttt{dynamic\_tanh}, \texttt{derf}, \texttt{rms\_norm}, \texttt{prms\_norm}, \texttt{flash\_norm}\}
\]
Las formulaciones más relevantes para este código base son:

\subsubsection{LayerNorm (Línea Base)}
LayerNorm \citep{ba2016layer} es el normalizador por defecto; centra y reescala cada token por su media y varianza por token:
\[
\mu = \frac{1}{d}\sum_{i=1}^{d} x_i, \quad \sigma^2 = \frac{1}{d}\sum_{i=1}^{d}(x_i - \mu)^2, \quad y_i = \gamma_i \frac{x_i - \mu}{\sqrt{\sigma^2 + \epsilon}} + \beta_i.
\]

\subsubsection{RMSNorm}
RMSNorm elimina la centralización de la media y solo reescala por la magnitud de la raíz cuadrada media \citep{zhang_root_2019}:
\[
\mathrm{RMS}(x)=\sqrt{\frac{1}{d}\sum_{i=1}^{d}x_i^2+\epsilon},\qquad
y_i=\gamma_i\frac{x_i}{\mathrm{RMS}(x)}.
\]
En comparación con LayerNorm, RMSNorm es computacionalmente más simple (sin resta de la media de características) y se usa a menudo cuando es importante reducir la sobrecarga de normalización (p.\,ej., Llama, GPT-NeoX). Tiene una escala $\gamma$ aprendible por dimensión pero sin término de bias.

\subsubsection{RMSNorm Parcial (pRMSNorm)}
pRMSNorm \citep{zhang_root_2019} reduce aún más la sobrecarga estimando el RMS de solo el primer $p$\% de las dimensiones ocultas:
\[
\mathrm{RMS}(x)=\sqrt{\frac{1}{k}\sum_{i=1}^{k}x_i^2+\epsilon},\qquad k=\lceil d\cdot p\rceil,
\]
explotando el supuesto de que las neuronas dentro de una capa son aproximadamente i.i.d. La misma escala $\gamma$ se aplica a todas las dimensiones. La razón parcial $p$ se controla mediante el campo de configuración \texttt{prms\_partial\_ratio} (por defecto $6.25\%$). El paper reporta que se pueden lograr resultados competitivos con tan solo el $6.25\%$ de las entradas, aunque puede surgir inestabilidad de gradiente con razones muy pequeñas.

\subsubsection{Tanh Dinámico (DyT)}
El Tanh Dinámico propone reemplazar la normalización explícita con un mapeo elemento a elemento acotado \citep{zhu_transformers_2025}:
\[
\mathrm{DyT}(x)=\tanh(\alpha x)
\]
donde $\alpha$ se aprende. La idea central es que la contracción no lineal acotada puede proporcionar un escalado de señal estable sin calcular explícitamente estadísticas de normalización por token.

\subsubsection{Erf Dinámico (Derf)}
Derf extiende la misma dirección libre de normalización usando un mapeo basado en la función error \citep{chen_stronger_2025}:
\[
\mathrm{Derf}(x)=\mathrm{erf}(\alpha x+s)
\]
con escala/desplazamiento aprendibles. Los resultados reportados en el trabajo citado indican un rendimiento superior al de DyT y las líneas base de normalización comunes en múltiples dominios.

\subsubsection{FlashNorm (RMSNorm sin pesos)}
FlashNorm \citep{graef2026flashnorm} no es una nueva normalización matemática; es una \emph{reescritura algebraica exacta} del ubicuo par ``$\mathrm{RMSNorm} \to \mathrm{Lineal}$'' que (i) elimina la escala aprendible por dimensión $\mathbf{g}$ al plegarla en los pesos de la capa lineal subsecuente, y (ii) difiere la reducción escalar RMS a la salida de la multiplicación matricial, de modo que el \texttt{matmul} (unidad matricial) y la reducción RMS (unidad vectorial) se ejecuten en paralelo. Tres proposiciones formalizan la reescritura.

\paragraph{Proposición 1 (Plegado de pesos).}
Dada la activación $\mathbf{a}\in\mathbb{R}^{n}$, el peso RMSNorm $\mathbf{g}\in\mathbb{R}^{n}$ y una lineal subsecuente $\mathbf{W}\in\mathbb{R}^{n\times k}$,
\[
\mathbf{z}=\mathrm{RMSNorm}(\mathbf{a})\mathbf{W}=\frac{\mathbf{a}}{\mathrm{RMS}(\mathbf{a})}\mathbf{W}^{*},\qquad
W^{*}_{i,j}=g_{i}\cdot W_{i,j}\ \Leftrightarrow\ \mathbf{W}^{*}=\operatorname{diag}(\mathbf{g})\mathbf{W}.
\]
La escala por dimensión $\mathbf{g}$ se vuelve redundante y se pliega en $\mathbf{W}$ una única vez al cargar el modelo. El módulo autónomo \texttt{FlashNorm} de este repositorio aplica esto directamente: \emph{no tiene parámetros aprendibles} (la escala está destinada a ser absorbida por la proyección subsecuente).

\paragraph{Proposición 2 (Normalización diferida).}
Para una lineal sin bias $\mathbf{W}^{*}$, la conmutatividad escalar--matricial produce
\[
\frac{\mathbf{a}}{\mathrm{RMS}(\mathbf{a})}\mathbf{W}^{*}=\left(\mathbf{a}\mathbf{W}^{*}\right)\cdot\frac{1}{\mathrm{RMS}(\mathbf{a})}.
\]
El \texttt{matmul} $\mathbf{a}\mathbf{W}^{*}$ y la reducción RMS son independientes entre sí; en hardware con unidades matriciales y vectoriales diferenciadas se ejecutan en paralelo, y solo una multiplicación vectorial--escalar los sincroniza. El módulo fusionado \texttt{FlashNormLinear} implementa esta reescritura: en la ruta sin bias calcula el \texttt{matmul} sobre la entrada \emph{no normalizada} y multiplica el resultado por el RMS recíproco por token. Cuando la lineal tiene bias $\mathbf{c}$, la reescritura deja de ser exacta (Observación 1 del paper): $\left(\mathbf{a}\mathbf{W}^{*}+\mathbf{c}\right)/\mathrm{RMS}(\mathbf{a})\neq\mathbf{a}\mathbf{W}^{*}/\mathrm{RMS}(\mathbf{a})+\mathbf{c}$; la capa entonces aplica el escalar a la entrada primero y deja que la lineal añada su bias como de costumbre.

\paragraph{Proposición 3 (Cancelación de RMS).}
Por la invariancia de escala de RMS, $\mathrm{RMS}(\alpha\mathbf{a})=\alpha\,\mathrm{RMS}(\mathbf{a})$. Por tanto, un patrón ``$\mathrm{RMSNorm}\to\mathrm{Lineal}\to\mathrm{RMSNorm}$'' permite que el \emph{primer} RMSNorm se elimine por completo:
\[
\mathrm{RMSNorm}\!\left(\frac{\mathbf{a}}{\mathrm{RMS}(\mathbf{a})}\mathbf{W}^{*}\right)=\mathrm{RMSNorm}\!\left(\mathbf{a}\mathbf{W}^{*}\right).
\]
Esta cancelación es relevante para la normalización QKV (Gemma~4) y la normalización latente MLA (DeepSeek-V2, Mistral Small~4); no se aplica automáticamente en este repositorio pero se documenta por completitud.

\paragraph{Composición con RMS parcial.}
FlashNorm se compone naturalmente con el truco pRMSNorm (§\ref{sec:annex-normalization}, subsección pRMSNorm): cuando \texttt{flashnorm\_partial\_ratio} $=p>0$, el RMS se estima solo de los primeros $k=\lceil d\cdot p\rceil$ elementos de la entrada, reduciendo aún más el costo de la reducción RMS en la unidad vectorial.

\paragraph{Interacción con BitNet.}
Plegar un $\mathbf{g}$ de coma flotante en el tensor de pesos ternarios de BitNet $\{-1,0,+1\}$ \citep{wang_bitnet_2023} destruiría la cuantización, y BitLinear aplica su propio LayerNorm interno antes de la cuantización de activaciones, el cual necesita la entrada normalizada. El módulo fusionado \texttt{FlashNormBitLinear} por tanto recurre a la composición secuencial $\mathrm{FlashNorm}\to\mathrm{BitLinear}$: el \texttt{matmul} y la reducción RMS aún se ejecutan en paralelo, pero la multiplicación escalar post-\texttt{matmul} de la Prop.~2 se sacrifica. El plegado de $\mathbf{g}$ en coma flotante en un tensor de pesos de 1.58 bits se deja a trabajo futuro a nivel kernel.

\paragraph{Pseudocódigo PyTorch.}
\begin{algorithmic}[1]
\State \textbf{Entrada:} activación $\mathbf{a}\in\mathbb{R}^{T\times n}$, peso RMSNorm $\mathbf{g}\in\mathbb{R}^{n}$, lineal $\mathbf{W}\in\mathbb{R}^{n\times k}$, bias $\mathbf{c}\in\mathbb{R}^{k}$ o \texttt{None}
\State \texttt{// Prop. 1: plegado de pesos al cargar}
\State $\mathbf{W}^{*}\gets \mathbf{g}\odot\mathbf{W}$ \Comment{por columnas; \texttt{g.unsqueeze(0) * W} en PyTorch}
\State $\mathbf{g}\gets \mathbf{1}$ \Comment{la norma autónoma ahora es sin pesos}
\State \texttt{// Prop. 2: paso forward con RMS diferida (ruta sin bias)}
\State $\mathbf{b}\gets \mathbf{a}\mathbf{W}^{*\top}$ \Comment{unidad matricial / tensor cores}
\State $\mathrm{rms\_inv}\gets \mathrm{rsqrt}(\mathrm{mean}(\mathbf{a}^{2})+\varepsilon)$ \Comment{unidad vectorial, independiente de $\mathbf{b}$}
\If{$\mathbf{c}=\texttt{None}$}
    \State \textbf{retornar} $\mathbf{b}\cdot\mathrm{rms\_inv}$ \Comment{una sola multiplicación vectorial--escalar}
\Else
    \State \texttt{// Obs. 1: ruta con bias secuencial; matmul/RMS paralelizan}
    \State \textbf{retornar} $(\mathbf{a}\cdot\mathrm{rms\_inv})\mathbf{W}^{*\top}+\mathbf{c}$
\EndIf
\end{algorithmic}

\paragraph{Latencia reportada (GPU NVIDIA T4).}
El paper reporta las siguientes aceleraciones sobre la operación \emph{norm-then-project}, con un despacho adaptativo que cae al camino secuencial en el régión limitada por memoria (decodificación):

\scriptsize
\begin{tabular}{llrrrl}
\toprule
\textbf{Escala} & \textbf{Tokens} & \textbf{Secuencial (ms)} & \textbf{FlashNorm (ms)} & \textbf{Aceleración} & \textbf{Notas} \\
\midrule
SmolLM2-135M & 4096 & 0.704 & 0.468 & $+33.6\%$ & limitada por cómputo (prefill) \\
SmolLM2-135M & 8192 & 0.929 & 0.599 & $+35.5\%$ & limitada por cómputo (prefill) \\
Llama-7B     & 1024 & 1.882 & 1.654 & $+12.1\%$ & limitada por cómputo (prefill) \\
Llama-7B     & 4096 & 7.628 & 6.570 & $+13.9\%$ & limitada por cómputo (prefill) \\
\bottomrule
\end{tabular}
\normalsize

\paragraph{Validación de plegado sin pérdida.}
El plegado de pesos (Prop.~1) es matemáticamente exacto; el paper verificó salidas bit-idénticas en tres checkpoints abiertos:

\scriptsize
\begin{tabular}{lll}
\toprule
\textbf{Modelo} & \textbf{Precisión} & \textbf{Resultado} \\
\midrule
SmolLM2-135M & fp16 & diff máxima $=0.0$, similitud coseno $=1.0$, texto idéntico \\
Llama-3.2-1B & fp32 & generación voraz bit-idéntica \\
Llama-3.1-8B & fp32 & generación voraz bit-idéntica \\
\bottomrule
\end{tabular}
\normalsize
La preservación de calidad en \texttt{lm-evaluation-harness} (wikitext word\_ppl, MMLU 5-shot, HellaSwag 0-shot) se mantiene dentro del ruido del benchmark (peor $\Delta=+0.0017$ word\_ppl).

\subsubsection{Comparación}
\scriptsize
\begin{longtable}{L{2.2cm}L{3.4cm}L{2.0cm}L{4.0cm}}
\toprule
\textbf{Método} & \textbf{Fórmula} & \textbf{Estadísticas Necesarias} & \textbf{Notas} \\
\midrule
\endhead
LayerNorm & $\gamma_i (x_i - \mu)/\sqrt{\sigma^2 + \epsilon} + \beta_i$ & media + varianza & Por defecto; centrado y reescalado completo \citep{ba2016layer}. \\
RMSNorm & $\gamma_i x_i/\sqrt{\frac{1}{d}\sum_j x_j^2+\epsilon}$ & Solo RMS & Menor sobrecarga; línea base ampliamente usada \citep{zhang_root_2019}. \\
pRMSNorm & $\gamma_i x_i/\sqrt{\frac{1}{k}\sum_{j\le k} x_j^2+\epsilon}$ & RMS parcial & RMS del primer $p$\% de dims; \texttt{prms\_partial\_ratio} por defecto $6.25\%$ \citep{zhang_root_2019}. \\
Tanh Dinámico & $\tanh(\alpha x)$ & ninguna & Transformación acotada libre de normalización; reemplazo directo \citep{zhu_transformers_2025}. \\
Erf Dinámico & $\mathrm{erf}(\alpha x+s)$ & ninguna & Alternativa libre de normalización; mejora sobre DyT \citep{chen_stronger_2025}. \\
FlashNorm & $x_i/\sqrt{\frac{1}{d}\sum_j x_j^2+\epsilon}$ (sin pesos) & Solo RMS & Reescritura exacta de RMSNorm$\to$Lineal: pliega $\mathbf{g}$ en $\mathbf{W}$ (Prop.~1), difiere el escalar RMS a la salida del \texttt{matmul} (Prop.~2); 12--35\% de reducción de latencia en GPU T4 \citep{graef2026flashnorm}. \\
\bottomrule
\end{longtable}
\normalsize

\subsection{Normalización Query--Key y Control de Logits de Atención}
\label{sec:annex-qk-norm-es}

Los logits de atención $\ell_{i,t,h} = \mathbf{q}_{i,h}^\top \mathbf{k}_{t,g(h)} / \sqrt{d_h}$ son la entrada del softmax. Su crecimiento descontrolado es una fuente clásica de inestabilidad de entrenamiento: saturación de la atención, colapso de entropía, picos de pérdida y, a escala, divergencia \citep{henry2020querykey,kimi2025k2}. Este repositorio expone cinco mecanismos bajo \texttt{norm.qk\_norm.method}; tres actúan en el forward y dos sobre el estado del optimizador. Son mutuamente excluyentes por configuración.

\paragraph{Normalización QK en el forward.}
La normalización QK \citep{henry2020querykey} normaliza queries y keys a lo largo de la dimensión de características por cabeza antes del producto punto. En forma RMSNorm,
\[
\hat{\mathbf{q}}_{i,h} = \frac{\mathbf{q}'_{i,h}\odot\boldsymbol{\gamma}_q}{S_q(\mathbf{q}'_{i,h})},\qquad
S_q(\mathbf{q}) = \sqrt{\frac{1}{d_h}\sum_{j=1}^{d_h} q_j^2 + \epsilon},
\]
e idénticamente para $\hat{\mathbf{k}}$ con $\boldsymbol{\gamma}_k$. Los pesos afines se comparten entre cabezas y se difunden por el eje de cabezas, por lo que el coste paramétrico es $2d_h$ (\texttt{rms\_norm}) o $2\cdot 2d_h$ con sesgo (\texttt{layer\_norm}); \texttt{l2\_norm} no tiene pesos. Dado que las normas normalizadas cumplen $\|\hat{\mathbf{q}}\|_2 \approx \|\boldsymbol{\gamma}_q\|_2$ salvo dirección y $\epsilon$, el crecimiento residual de los logits queda mediado por el producto afín aprendido $\boldsymbol{\gamma}_q\odot\boldsymbol{\gamma}_k$ en lugar de por normas crudas de query/key sin control.

\paragraph{Normalización QK con Multi-Head Latent Attention.}
MLA cachea un latente de bajo rango $\mathbf{c}_t$ y reconstruye las keys por up-proyección, por lo que la QK RMSNorm post-proyección parecería requerir las keys completas que MLA evita cachear. \citet{han2026qknormed} muestran que la incompatibilidad es un artefacto de implementación: RMSNorm se descompone en un peso afín estático y una estadística escalar dinámica,
\[
\hat{\mathbf{k}}^c_{t,g}=\frac{(\mathbf{c}_t\mathbf{W}^K_g)\odot\boldsymbol{\gamma}^c_k}{S^c_k(\mathbf{c}_t\mathbf{W}^K_g)},\qquad
\ell^c_{i,t,h}=\big((\mathbf{q}'^{c}_{i,h}\odot\boldsymbol{\gamma}^c_q\odot\boldsymbol{\gamma}^c_k)(\mathbf{W}^K_{g(h)})^\top\big)\cdot\mathbf{c}_t\;\alpha^q_{i,h}\alpha^k_{t,g(h)},
\]
donde $\alpha^q=1/S_q$ y $\alpha^k=1/S_k$ son escalares. El peso estático se absorbe en la proyección del lado de la query y solo se cachea un escalar de RMS inverso por token y grupo KV ($T\times G$ escalares frente a $T\times G\times d_c$ keys de contenido completas), preservando la ruta latente de decodificación; el paper reporta menor pérdida que QK-Clip y $<2\%$ de sobrecarga de decodificación hasta 256k tokens. En este repositorio la variante MLA de \citet{mehta_mla_2025} ya re-expande las keys desde el latente cacheado en cada paso, de modo que la normalización explícita por cabeza de $\mathbf{q}$ y $\mathbf{k}$ descomprimidos es \emph{exactamente} equivalente a la formulación absorbida tanto en entrenamiento como en decodificación con caché; no se necesita un caché escalar. Los bloques de contenido y RoPE se normalizan por bloques, siguiendo la frontera natural entre la ruta latente de contenido y la ruta rotatoria materializada.

\paragraph{QK-Clip.}
MuonClip de Kimi K2 \citep{kimi2025k2} ataca el mismo modo de fallo después del producto punto. Para cada cabeza $h$ define el logit pre-softmax máximo por cabeza sobre el lote,
\[
S^h_{\max}=\frac{1}{\sqrt{d}}\max_{\mathbf{X}\in\mathcal{B}}\max_{i,j}\mathbf{Q}^h_i{\mathbf{K}^h_j}^\top,
\]
y, tras cada paso del optimizador, reescala los pesos de proyección con
\[
\gamma_h=\min\!\left(1,\frac{\tau}{S^h_{\max}}\right),\qquad
\mathbf{W}^h_q\leftarrow\gamma_h^{\alpha}\mathbf{W}^h_q,\qquad
\mathbf{W}^h_k\leftarrow\gamma_h^{1-\alpha}\mathbf{W}^h_k .
\]
La intervención nunca altera el forward/backward del paso actual; las cabezas ya por debajo de $\tau$ reciben $\gamma_h=1$, por lo que el mecanismo se autodesactiva cuando el entrenamiento se estabiliza (K2 usó $\tau=100$ y observó todas las cabezas por debajo del umbral tras aproximadamente el 30\% de los pasos). Se aplica por cabeza a los bloques de filas de \texttt{q\_proj} y \texttt{k\_proj}; para atención latente el lado de la key es \texttt{uk\_proj}, y para grouped-query attention cada cabeza KV recibe el menor $\gamma_h$ de las cabezas de query que la leen, lo que acota todos los logits del grupo. \texttt{clip\_alpha} vale $0.5$ por defecto (es decir, $\sqrt{\gamma_h}$ en cada lado, la elección balanceada de K2).

\paragraph{QuacK.}
En lugar de una intervención en el forward o en los pesos, QuacK \citep{anson2025controlling} asigna learning rates dependientes del parámetro a los pesos de query y key. Guardando las normas de Frobenius en la inicialización, antes de cada paso fija
\[
\eta_{W_Q}=\tau\,\eta\,\frac{\|\mathbf{W}_K\|_{\mathrm{init}}}{\|\mathbf{W}_K\|},\qquad
\eta_{W_K}=\tau\,\eta\,\frac{\|\mathbf{W}_Q\|_{\mathrm{init}}}{\|\mathbf{W}_Q\|},
\]
lo que se sigue de controlar el cambio de primer orden $\Delta\ell\approx(\mathbf{q}\,\Delta\mathbf{k}^\top+\Delta\mathbf{q}\,\mathbf{k}^\top)/\sqrt{d}$: cuando un lado crece, el learning rate del otro decrece. Las normas de Frobenius y espectral se comportaron de forma indistinguible en el paper, por lo que se usa Frobenius. La tasa relativa $\tau$ vale $1.0$ por defecto; el paper reporta una estabilidad comparable a QK-norm con learning rates base altos y alrededor de un 10\% de aceleración porque desaparecen dos reducciones RMS por bloque. Aquí el mecanismo es independiente del optimizador: el trainer separa los pesos de query/key en grupos de parámetros dedicados, de modo que cualquiera de los optimizadores registrados puede soportarlo.

\paragraph{Contrato de esquema y validación.}
\texttt{norm.qk\_norm} es un sub-objeto con \texttt{enabled}, \texttt{method} $\in$ \{\texttt{rms\_norm}, \texttt{layer\_norm}, \texttt{l2\_norm}, \texttt{qk\_clip}, \texttt{quack}\}, \texttt{clip\_tau}, \texttt{clip\_alpha}, \texttt{quack\_tau} y \texttt{eps}. El mecanismo requiere proyecciones explícitas de query/key por cabeza: activarlo para un patrón de capas que contenga un mixer recurrente o de espacio de estados se rechaza al cargar (\texttt{QK\_NORM\_MIXERS}); \texttt{qk\_clip} requiere además logits de producto punto por cabeza y por tanto excluye el router de mezcla Gaussiana (\texttt{QK\_CLIP\_MIXERS}).

\paragraph{Pseudocódigo PyTorch (QK-Clip).}
\begin{algorithmic}[1]
\State \textbf{Entrada:} modelo, umbral $\tau$, exponente $\alpha$, $S^h_{\max}$ por cabeza cacheado
\For{cada cabeza de atención $h$ con $S^h_{\max}>\tau$}
    \State $\gamma_h\gets\tau/S^h_{\max}$
    \State $\mathbf{W}^h_q\gets\gamma_h^{\alpha}\mathbf{W}^h_q$ \Comment{bloque de filas $[h d_h:(h+1)d_h]$}
    \State $\mathbf{W}^h_k\gets\gamma_h^{1-\alpha}\mathbf{W}^h_k$ \Comment{\texttt{k\_proj} o \texttt{uk\_proj}; mínimo sobre el grupo GQA}
\EndFor
\State limpiar $S^h_{\max}$ \Comment{evitar reaplicar una estadística obsoleta}
\end{algorithmic}

\paragraph{Comparación.}
\scriptsize
\begin{longtable}{L{2.1cm}L{3.5cm}L{2.4cm}L{4.3cm}}
\toprule
\textbf{Método} & \textbf{Intervención} & \textbf{Estado extra} & \textbf{Notas} \\
\midrule
\endhead
\texttt{rms\_norm} & RMSNorm por cabeza en el forward & $2d_h$ & QK-norm estándar; exacta con el MLA re-expansivo de este repo \citep{henry2020querykey,han2026qknormed}. \\
\texttt{layer\_norm} & LayerNorm por cabeza en el forward & $4d_h$ & Variante con centrado de media (SAN); algo más costosa. \\
\texttt{l2\_norm} & L2 por cabeza en el forward & ninguno & Ablación sin pesos. \\
\texttt{qk\_clip} & reescalado de pesos post-paso & $S_{\max}$ por cabeza & Se autodesactiva; necesita logits por cabeza \citep{kimi2025k2}. \\
\texttt{quack} & learning rates de query/key por paso & normas iniciales & Independiente del optimizador; $\sim$10\% más barato que QK-norm \citep{anson2025controlling}. \\
\bottomrule
\end{longtable}
\normalsize

